% ===== 原文 pp.99--113 =====
% [原文 p.99]
\section{第二讲（Whitehead 群与 Steinberg 群）}

我们用 $\GL_nA$ 记元素取自环 $A$ 的 $n\times n$ 可逆矩阵的集合（实际上它是乘法群）。显然
\[
  \GL_nA \subset \GL_{n+1}A \subset \dots \subset \GL A = \cup_n \GL_nA,
\]
其中嵌入由
\[
  a \into \begin{pmatrix} a & 0 \\ 0 & 1 \end{pmatrix}
\]
给出。
% [Tutor Note: p.99 原文此处箭头印作 $\hookrightarrow$（1400 dpi 复核：$a\hookrightarrow\begin{pmatrix}a&0\\0&1\end{pmatrix}$），按文意应为 $\mapsto$（$a$ 映到分块矩阵）；译文照页图作 $\into$。]

我们用 $e_{ij}$ 记在 $(i,j)$ 位置为 $1$、其余位置全为 $0$ 的矩阵（$i\ne j$）。显然
\[
  e_{ij}e_{kl} = \delta_{jk}e_{il}.
\]
设 $a\in A$。则我们令
\[
  e_{ij}^a = 1 + ae_{ij},
\]
而
\[
  e_{ij}^a e_{ij}^b = e_{ij}^{a+b}.
\]
% [Tutor Note: p.99 这里的 $e_{ij}^a$ 不是矩阵 $e_{ij}$ 的 $a$ 次幂，而是初等矩阵 $1+ae_{ij}$；$e_{ij}$ 是矩阵单位（$(i,j)$ 处为 $1$、其余为 $0$），故 $e_{ij}e_{kl}=\delta_{jk}e_{il}$ 只在 $j=k$ 时不等于零。核对页图：上标 $a$ 是环元素，下标 $i,j$ 才是位置。]
回忆在乘法群中，两个元素的换位子（commutator）定义为
\[
  [x,y] = xyx^{-1}y^{-1},
\]
而反换位子定义为
\[
  [x,y]^{-1} = [y,x].
\]
利用上述定义，我们计算换位子
\[
  [e_{ij}^a, e_{kl}^b] =
  \begin{cases}
    1, & \text{若 } j\ne k,\ i\ne l \\
    e_{il}^{ab}, & \text{若 } j=k,\ i\ne l \\
    e_{kj}^{-ba}, & \text{若 } j\ne k,\ i=l
  \end{cases}
\]
此外我们还计算下列量：
\[
  e_{ij}^a e_{jk}^b = (1+ae_{ij})(1+be_{jk}) = 1 + ae_{ij} + be_{jk} + abe_{ik},
\]

% [原文 p.100]
\[
  e_{ij}^a e_{jk}^b e_{ij}^{-a} = (1+ae_{ij}+be_{jk}+abe_{ik})(1-ae_{ij}) = 1 + ae_{ij} + be_{jk} + abe_{ik} - ae_{ij},
\]
\[
  [e_{ij}^a, e_{jk}^b] = (1+ae_{ij}+be_{jk}+abe_{ik})(1-be_{jk}) = 1 + be_{jk} + abe_{ik} - be_{jk} = e_{ik}^{ab}
\]
\[
  [e_{ij}^a, e_{ki}^b]^{-1} = [e_{ki}^b, e_{ij}^a] = e_{kj}^{ba} \Rightarrow [e_{ij}^a, e_{ki}^b] = e_{kj}^{-ba}.
\]

\begin{defn}{1}
我们用 $E_nA$ 记 $\GL_nA$ 中由元素 $e_{ij}^a$（$1\le i,j\le n$，$i\ne j$，$a\in A$）生成的子群。此外我们用 $EA$ 记所有 $E_nA$ 的并：
\[
  EA = \cup_n E_nA.
\]
\end{defn}

\begin{defn}{2}
群 $G$ 称为\emph{完美群}（perfect group），如果它等于其换位子子群
\[
  G = [G,G],
\]
其中 $[G,G]$ 是由所有形如 $[g,g']$（$\forall g,g'\in G$）的元素生成的 $G$ 的子群。
\end{defn}

\begin{defn}{3}
我们如下定义 $G$ 的极大阿贝尔商群 $G_{(ab)}$：
\[
  G_{(ab)} = G/[G,G].
\]
（$ab$ 代表「abelian」，它们不是指标。）
\end{defn}

\begin{prop}{1}
群 $E_nA$ 对 $n\ge 3$ 是完美的。
\end{prop}

% [Tutor Note: p.100 证明中的取法：对任意给定的 $i,k$，只要 $n\ge 3$ 就能取到 $j\ne i$ 且 $j\ne k$，于是每个生成元 $e_{ik}^a$ 都是换位子；这正是「$n\ge 3$」不可省的缘由。$e_{jk}^1$ 的上标是 $1\in A$，不是单位矩阵。]
\begin{proof}
由于
\[
  e_{ik}^a = [e_{ij}^a, e_{jk}^1] \in [E_nA, E_nA],
\]
给定 $i,k$，我们选取 $j\ne i$ 且 $j\ne k$，便看出所有生成元都是换位子。
\end{proof}

\begin{lem}{1}（Whitehead 引理）
\[
  EA = [\GL A,\GL A] = [EA,EA].
\]
\end{lem}

% [Tutor Note: p.100--101 证明的关键是 $2n\times 2n$ 分块技巧：把 $a\in\GL_nA$ 放在左上块，则 $\begin{pmatrix}1&a\\0&1\end{pmatrix}$ 是形如 $e_{ij}^{a_{ij}}$（$1\le i\le n$，$n+1\le j\le 2n$）的初等矩阵之积，故属于 $E_{2n}A$；再用 $\begin{pmatrix}0&a\\-a^{-1}&0\end{pmatrix}$ 把对角块 $\operatorname{diag}(a,a^{-1})$ 也化归为初等矩阵。]
\begin{proof}
设 $a\in\GL_nA$，于是
\[
  \begin{pmatrix} 1 & a \\ 0 & 1 \end{pmatrix} = \prod e_{ij}^{a_{ij}} \in \GL_{2n}A
\]
% [原文 p.101]
其中 $1\le i\le n$ 且 $n+1\le j\le 2n$。则
\[
  \begin{pmatrix} 1 & a \\ 0 & 1 \end{pmatrix}
  \begin{pmatrix} 1 & 0 \\ -a^{-1} & 1 \end{pmatrix}
  \begin{pmatrix} 1 & a \\ 0 & 1 \end{pmatrix}
  = \begin{pmatrix} 0 & a \\ -a^{-1} & 0 \end{pmatrix} \in \GL_{2n}A.
\]
此外
\[
  \begin{pmatrix} a & 0 \\ 0 & a^{-1} \end{pmatrix}
  = \begin{pmatrix} 0 & a \\ -a^{-1} & 0 \end{pmatrix}
  \begin{pmatrix} 0 & -1 \\ 1 & 0 \end{pmatrix} \in E_{2n}A.
\]
接下来，设 $a,b\in\GL_nA$。则对于上述（分块形式的）对角矩阵，我们有
\[
  \resizebox{0.98\textwidth}{!}{$\displaystyle
  \operatorname{diag}([a,b],1,1) = \operatorname{diag}(a,a^{-1},1)\operatorname{diag}(b,1,b^{-1})[\operatorname{diag}(a^{-1},a^{-1},1)]^{-1}[\operatorname{diag}(b,1,b^{-1})]^{-1}$}
\]
（它们全都属于群 $E_{3n}A$）。因此 $[\GL_nA,\GL_nA]$ 在 $\GL_{3n}A$ 中包含于 $E_{3n}A$ 之中。对所有 $n$ 取并，我们发现 $[\GL_nA,\GL_nA]\subset EA$。因此
\[
  EA = [EA,EA] \subset [\GL A,\GL A] \subset EA.
\]
\end{proof}

% [Tutor Note: p.101 上式第三个因子原文作 $[\operatorname{diag}(a^{-1},a^{-1},1)]^{-1}$（500 dpi 复核无误），但它应为 $[\operatorname{diag}(a,a^{-1},1)]^{-1}$：令 $D_1=\operatorname{diag}(a,a^{-1},1)$、$D_2=\operatorname{diag}(b,1,b^{-1})$，则 $D_1D_2D_1^{-1}D_2^{-1}=\operatorname{diag}([a,b],1,1)$；而照原文所印公式算得的是 $\operatorname{diag}(abab^{-1},1,1)$（除 $b=1$ 外不等于 $\operatorname{diag}([a,b],1,1)$）。正文照译原文，此处指出正确形式。]

\begin{defn}{4}
环 $A$ 的第一个 K 群 $K_1A$ 定义为 $\GL A$ 的极大阿贝尔商群，即
\[
  K_1A = \GL A_{(ab)} = \GL A/EA.
\]
\end{defn}

第一个 K 群称为\emph{Whitehead 群}，以纪念定义了它的伟大的英国拓扑学家 John Whitehead。

\para{例}
\begin{enumerate}[label=\arabic*.]
\item 设 $A=\F$ 为域。用矩阵 $e_{ij}^a$ 左乘会把第 $j$ 行加到第 $i$ 行上。已知
\[
  E_n(\F) = \operatorname{ker}\{\GL_n(\F) \xra{\det} \F^x\}
\]
% [原文 p.102]
其中 $\F^x$ 表示在 $x$ 作用下的非零元素，上述映射由矩阵的行列式给出。接下来我们考虑商
\[
  \GL_n(\F)/E_n(\F) = \F^x = K_1(\F).
\]
\end{enumerate}

% [Tutor Note: p.102 按通常记号 $\F^x=\F\setminus\{0\}$ 是域的乘法群（$\F$ 的单位群）；「在 $x$ 作用下的非零元素」一句原文含混，应理解为 $\F^\times$ 即非零元乘法群。]
\begin{defn}{5}
Steinberg 群 $\operatorname{St}_nA$（$\operatorname{St}A$ 是 $n\to\infty$ 的情形）是以 $x_{ij}^a$（$i\ne j$，$a\in A$）为生成元的群，其关系为
\[
  x_{ij}^a x_{ij}^b = x_{ij}^{a+b}.
\]
换位子由下列关系给出：
\[
  [x_{ij}^a, x_{kl}^b] =
  \begin{cases}
    1, & \text{若 } j\ne k,\ i\ne l \\
    x_{ik}^{ab}, & \text{若 } j=k,\ i\ne l \\
    x_{kj}^{-ba}, & \text{若 } j\ne k,\ i=l
  \end{cases}
\]
\end{defn}

% [Tutor Note: p.102 定义 5 的抽象关系恰好是把 p.99--100 中 $e_{ij}^a$ 的乘法关系与换位子公式「照抄」成生成元之间的形式关系；正因如此才有典范满射 $\phi:\operatorname{St}A\to EA$，而 $\phi$ 的核正是 $K_2A$。注意：两处 cases 的三行条件相同，但第二行指标并不一致——p.99 的 $e$ 版作 $e_{il}^{ab}$，定义 5（及 p.104 的重述）却作 $x_{ik}^{ab}$；按换位子关系 $[x_{ij}^a,x_{jl}^b]=x_{il}^{ab}$，$x_{ik}$ 应为 $x_{il}$ 之误。译文两处均照页图。]
存在典范满射
\[
  \phi : \operatorname{St}A \to EA
\]
使得
\[
  \phi(x_{ij}^a) = e_{ij}^a.
\]

\begin{prop}{2}
上述 $\phi$ 的核是群 $\operatorname{St}A$ 的中心。
\end{prop}

\begin{defn}{6}
第二个 K 群 $K_2A$ 定义为上述典范满射 $\phi$ 的核：
\[
  K_2A = \operatorname{Ker}\{\phi : \operatorname{St}A \to EA\}.
\]
\end{defn}

\begin{proof}[命题 2 的证明]
设 $C_n$ 是 $\operatorname{St}A$ 中由 $x_{in}^a$（$i\ne n$，$a\in A$）生成的子群。设 $\phi(C_n)$ 是 $EA$ 中由 $e_{in}^a$（$i\ne n$，$a\in A$）定义的子群。

\para{断言} $\phi : C_n \to \phi(C_n)$ 是同构，即 $\phi$ 在 $C_n$ 上的限制是单射。

\emph{断言的证明：} 由于各个 $x_{in}^a$ 彼此交换，在由
\[
  (a_i)_{i\ne n} \mapsto \prod_{i\ne n} x_{in}^{a_i}
\]
定义的映射
\[
  \oplus_{i\ne n} A \to C_n
\]
% [原文 p.103]
中，乘积不依赖于次序，而且它是同态，因为：
\[
  (a_i+a_i') \mapsto \prod_{i\ne n} x_{in}^{a_i+a_i'} = \prod_{i\ne n} x_{in}^{a_i}x_{in}^{a_i'} = \prod_{i\ne n} x_{in}^{a_i}\prod_{i\ne n} x_{in}^{a_i'}.
\]
接下来我们定义群 $R_n$（它是群 $\operatorname{St}A$ 的子群），即由 $x_{ni}^a$（$i\ne n$，$a\in A$）生成的群。类似的论证证明 $\phi$ 在 $R_n$ 上的限制是单射。

为完成命题 2（即 $\phi$ 的核是 $\operatorname{St}A$ 的中心）的证明，我们考虑任意元素 $\alpha\in\operatorname{Ker}\phi$，并把它写成因子 $x_{ij}^a$ 的有限乘积。选取某个与 $\alpha$ 的这一具体表示中出现的所有 $i,j$ 都不同的 $n$。则 $\alpha$ 正规化 $C_n$，即
\[
  \alpha C_n\alpha^{-1} \subset C_n
\]
而
\[
  x_{ij}^a x_{kn}^b x_{ij}^{-a} =
  \begin{cases}
    x_{kn}^b, & k\ne i,j, \\
    x_{in}^{ab}x_{jn}^{-b}, & k=j
  \end{cases}
\]
以及
\[
  x_{ij}^a x_{jn}^b x_{ij}^{-a}x_{jn}^{-b} = x_{in}^{ab}.
\]
设 $\gamma\in C_n$，则 $\alpha\gamma\alpha^{-1}\in C_n$，并且
\[
  \phi(\alpha\gamma\alpha^{-1}) = \phi(\gamma)
\]
因为 $\alpha\in\operatorname{Ker}\phi$。但 $\phi$ 在 $C_n$ 上的限制是单射，从而 $\alpha\gamma\alpha^{-1}=\gamma$，于是 $\alpha$ 中心化 $C_n$。类似地，$\alpha$ 中心化 $R_n$。但并集 $C_n\cup R_n$ 生成 $\operatorname{St}A$，且中心为 1。
\end{proof}

% [Tutor Note: p.103 末句「the centre is 1」语义含混（照译）。作者此处实际只证了：$\operatorname{Ker}\phi$ 中的元中心化 $C_n$ 与 $R_n$，而 $C_n\cup R_n$ 生成 $\operatorname{St}A$，故 $\operatorname{Ker}\phi\subset Z(\operatorname{St}A)$；命题 2 所需的另一方向 $Z(\operatorname{St}A)\subset\operatorname{Ker}\phi$ 在原文中并未证明（p.105 只是重述「$K_2A$ 含于中心」，并未补证）。]

% [原文 p.104]
\section{第三讲（群的中心扩张）}

我们看到
\[
  K_1A = \GL A/EA.
\]
对拓扑空间 $X$ 有
\[
  K^{-1}(X) = [X,\GL\C],
\]
其中右端（如我们从同伦论中所知）表示从 $X$ 到 $\GL\C$ 的连续映射的同伦类集合，并且
\[
  \Hom_{(\text{spaces})}(X,\GL_n\C) = \GL_n(C(X)),
\]
其中 $C(X)$ 表示从 $X$ 到 $\C$ 的连续映射的集合（含幺环）。

% [Tutor Note: p.104 这里 $\GL\C$ 是直极限 $\varinjlim_n\GL_n\C$；$[X,\GL\C]$ 表示同伦类集合而 $\Hom_{(\text{spaces})}$ 表示连续映射空间，两者的区别正是「同伦类」与「映射空间」的区别（$\pi_0$ 之后二者一致）。]
设
\[
  \begin{pmatrix} 1 & ta \\ * & 1 \end{pmatrix}
\]
是从
\[
  \begin{pmatrix} 1 & 0 \\ 0 & 1 \end{pmatrix}
\]
到
\[
  \begin{pmatrix} 1 & a \\ * & 1 \end{pmatrix}
\]
的同伦。

设 $X$ 是拓扑空间，$SX$ 是它的悬挂。我们知道
\[
  K^{-2}(X) = K^{-1}(SX) = [SX,\GL\C] = [S^1,[X,\GL\C]]
\]
其中
\[
  [X,\GL\C] = \Hom(X,\GL\C) = \GL(C(X)).
\]
群 $\operatorname{St}A$ 有生成元 $x_{ij}^a$，$1\le i,j<\infty$，$a\in A$，以及关系
\[
  x_{ij}^a x_{ij}^b = x_{ij}^{a+b}
\]
和换位子
\[
  [x_{ij}^a, x_{kl}^b] =
  \begin{cases}
    1, & \text{若 } j\ne k,\ i\ne l \\
    x_{ik}^{ab}, & \text{若 } j=k,\ i\ne l \\
    x_{kj}^{-ba}, & \text{若 } j\ne k,\ i=l
  \end{cases}
\]

% [原文 p.105]
有下述中心扩张
\[
  1 \to K_2A \to \operatorname{St}A \to EA \to 1,
\]
即 $K_2A$ 包含于 $\operatorname{St}A$ 的中心，其中
\[
  x_{ij}^a \mapsto e_{ij}^a.
\]

\begin{thm}{1}
给定中心扩张
\[
  1 \to C \to Y \xra{\phi} \operatorname{St}A \to 1,
\]
它分裂，即存在同态 $s : \operatorname{St}A \to Y$ 使得 $\phi\circ s = \id$。
\end{thm}

\begin{cor}{1}
若 $Y=[Y,Y]$，则 $Y$ 同构于 $\operatorname{St}A$。
\end{cor}

\para{定理 1 的证明}
基本想法如下：设 $y_1,y_2\in Y$ 满足 $\phi(y_1)=\phi(y_2)$，即 $y_1=cy_2$，其中 $c\in C$。则
\[
  [y_1,y'] = [cy_2,y'] = cy_2y'(cy_2)'y'^{-1} = cy_2y'y_2^{-1}c^{-1}y'^{-1} = [y_2,y'],
\]
即对 $x\in\operatorname{St}A$，$[\phi^{-1}(x),y']$ 是 $Y$ 中良定义的元素；类似地，$[\phi^{-1}(x),\phi^{-1}(x')]$ 也是 $Y$ 中良定义的元素。然后
\[
  x_{ij}^a = [x_{in}^a, x_{nj}^1]
\]
（$n\ne i,j$），所以我们试着用
\[
  s(x_{ij}^a) = [\phi^{-1}x_{in}^a, \phi^{-1}x_{nj}^1]
\]
来定义 $s$。困难之处在于证明它与 $n$ 的选取无关，下面我们将用 2 条引理来完成：

% [Tutor Note: p.105 上式中原稿写作 $(cy_2)'$（500 dpi 复核：确为一个撇号），应为 $(cy_2)^{-1}$；否则 $cy_2y'(cy_2)'y'^{-1}$ 无法化到 $cy_2y'y_2^{-1}c^{-1}y'^{-1}$。另注意 $\phi^{-1}(x)$ 表示原像（$\phi$ 不是单射，故并非映射的逆）。]
\begin{lem}{1}
\[
  [\phi^{-1}x_{ij}^a, \phi^{-1}x_{kl}^b] = 1
\]
若 $j\ne k$ 且 $i\ne l$。
\end{lem}

% [原文 p.106]
\begin{proof}
选取不同于 $k,l,i,j$ 的 $n$。则
\[
  [\phi^{-1}x_{kn}^b, \phi^{-1}x_{nl}^1] \subset \phi^{-1}[x_{kn}^b, x_{nl}^1] = \phi^{-1}x_{kl}^b.
\]
取 $v\in\phi^{-1}x_{kn}^b$ 与 $w\in\phi^{-1}x_{nl}^1$。再取 $u\in\phi^{-1}x_{ij}^a$，只需证明 $[u,[v,w]]=1$。有：
\[
  u[v,w]u^{-1} = [uvu^{-1}, uwu^{-1}] = [v,w]
\]
如果 $uvu^{-1}$ 与 $v$ 以及 $uwu^{-1}$ 与 $w$ 分别同余。但 $\phi(u),\phi(v)$ 交换，且 $\phi(u),\phi(w)$ 也交换。
\end{proof}

% [Tutor Note: p.106 「同余」（congruent）指 $uvu^{-1}$ 与 $v$ 属于 $\operatorname{Ker}\phi$ 的同一陪集，因而 $\phi(uvu^{-1})=\phi(v)$；这是全段推理的关键，原文只用一句话带过。]
回忆下列对任意群都成立的恒等式：
\[
  [x,[y,z]] = [xy,z][z,x][z,y]
\]
以及
\[
  [xy,z] = [x,[y,z]][y,z][x,z].
\]
利用它们我们得到第二条引理：

\begin{lem}{2}
设 $h,i,j,k$ 互不相同，$a,b,c\in A$。则
\[
  [\phi^{-1}x_{hi}^a, [\phi^{-1}x_{ij}^b, \phi^{-1}x_{jk}^c]] = [[\phi^{-1}x_{hi}^a, \phi^{-1}x_{ij}^b], \phi^{-1}x_{jk}^c].
\]
\end{lem}

\begin{proof}
取 $u\in\phi^{-1}x_{hi}^a$，$v\in\phi^{-1}x_{ij}^b$，$w\in\phi^{-1}x_{jk}^c$。由引理 1 我们有
\[
  [u,w] = 1
\]
而
\[
  [u,v] \subset \phi^{-1}[x_{hi}^a, x_{ij}^b] = \phi^{-1}x_{hj}^{ab}
\]
它与 $u,v$ 交换，而 $[v,w]$ 与 $v,w$ 交换。此外
\[
  [u,[v,w]] \subset [\phi^{-1}x_{hi}^a, \phi^{-1}x_{ik}^{bc}] \subset \phi^{-1}x_{ik}^{abc}
\]
它与 $u,v,w$ 交换。类似地，$[[u,v],w]$ 与 $u,v,w$ 交换。

还有：
\[
  [x,[y,z]] = [xy,z][z,x][z,y]
\]
% [原文 p.107]
以及
\[
  [u,[v,w]] = [uv,w][w,u][w,v] = [[u,v]vu,w][w,v] = [vu[u,v],w][w,v]
\]
而
\[
  [xy,z] = [x,[y,z]][y,z][x,z].
\]
于是
\[
  [u,[v,w]] = [uv,w][w,u][w,v] = [[u,v]vu,w][w,v] = [vu[u,v],w][w,v]
\]
等于
\[
  [u,[v,w]] = [uv,w][w,u][w,v] = [[u,v]vu,w][w,v] = [vu[u,v],w][w,v] =
\]
\[
  = [vu,[[u,v],w]][[u,v],w][vu,w][w,v] = [[u,v],w]
\]
而
\[
  s_{ij}^a := s(x_{ij}^a) = [\phi^{-1}x_{in}^a, \phi^{-1}x_{nj}^1].
\]
我们可以把引理 2 改写为：
\[
  [\phi^{-1}x_{hi}^a, \phi^{-1}x_{ik}^{bc}] = [\phi^{-1}x_{hj}^{ab}, \phi^{-1}x_{jk}^c]
\]
以及
\[
  [\phi^{-1}x_{hi}^a, \phi^{-1}x_{ik}^b] = [\phi^{-1}x_{hj}^a, \phi^{-1}x_{jk}^1],
\]
即它与 $n$ 无关。
\end{proof}

\begin{thm}{2}
每个完美群都有一个「泛」中心扩张：
\[
  \begin{tikzcd}
    1 \arrow[r] & H_2(G) \arrow[r] \arrow[d] & \wt G \arrow[r] \arrow[d] & G \arrow[r] \arrow[d, equal] & 1 \\
    1 \arrow[r] & C \arrow[r] & Y \arrow[r] & G \arrow[r] & 1
  \end{tikzcd}
\]
其中「泛」是指存在唯一的映射 $H_2(G) \to C$。
\end{thm}

% [Tutor Note: p.107 定理 2 的「泛」性质：对任意中心扩张 $1\to C\to Y\to G\to 1$（$G$ 完美），存在唯一的群同态 $H_2(G)\to C$ 使图表交换；这正是 $H_2(G)$ 作为 Schur 乘子分类中心扩张的精确含义。图中竖直双线表示恒等映射。]
% [原文 p.108]
前面的结果说明，群 $\operatorname{St}A$ 是群 $EA$ 的万有中心扩张。特别地
\[
  H_2(EA) = \operatorname{Ker}(\operatorname{St}A \to EA) := K_2A
\]
而
\[
  K_1A = \GL A/[\GL A,\GL A] = H_1(\GL A,\Z),
\]
\[
  K_2A = H_2(EA,\Z)
\]
以及
\[
  K_3A = H_3(\operatorname{St}A,\Z).
\]

% [Tutor Note: p.108 这里把 $K_1,K_2,K_3$ 解释为群同调：$K_1A=H_1(\GL A,\Z)=\GL A/[\GL A,\GL A]$（Hurewicz）；$K_2A=H_2(EA,\Z)$ 与 $K_2A=\operatorname{Ker}(\operatorname{St}A\to EA)$ 的一致正是上面的泛中心扩张定理；$K_3A=H_3(\operatorname{St}A,\Z)$ 是 Gersten 的定理，此处未给证明。核对：三个公式的下标与群名须一一对应。]

% [原文 p.109]
\section{第四讲（分类空间与群（上）同调）}

设 $G$ 是群，$M$ 是 $G$-模。我们知道对于上同调有
\[
  H^i(G,M) = H^i(C^{\cdot}(G,M))
\]
而同调也有类似的结果。若 $C$ 是带有平凡 $G$-作用的阿贝尔群，则我们有
\[
  H^0(G,C) = C,
\]
\[
  H^1(G,C) = \Hom(G,C) = \Hom(G_{(ab)},C)
\]
（其中 $G_{(ab)} = G/[G,G]$）以及
\[
  H^2(G,C) = \{\text{$G$ 被 $C$ 的中心扩张的同构类之集}\},
\]
即
\[
  \begin{tikzcd}
    1 \arrow[r] & C \arrow[r] \arrow[d, equal] & E \arrow[r] \arrow[d] & G \arrow[r] \arrow[d, equal] & 1 \\
    1 \arrow[r] & C \arrow[r] & E' \arrow[r] & G \arrow[r] & 1
  \end{tikzcd}
\]
其中 $C$ 是 $E$ 的中心的一个子集。

% [Tutor Note: p.109 $C^{\cdot}(G,M)$ 中的 $\cdot$ 是上标中心点（页图上确为抬高的点），表示上链复形；又 $H^i$ 是上同调（上标），不要与同调 $H_i$ 混写。]
群（上）同调的拓扑解释如下：$G$ 有一个分类空间 $BG$，它是一个带基点的好空间，使得：
\begin{enumerate}[label=\arabic*.]
\item $\pi_1(BG) = G$。
\item $BG$ 的万有覆盖空间是可缩的。
\end{enumerate}

例如若 $G=\Z$，则 $BG=S^1$。

$BG$ 在同伦等价意义下唯一。

\para{事实} $H^i(G,C) = H^i(BG,C)$。这对同调也成立。

由万有系数定理（见附录），可得
\[
  0 \to \Ext^1_{\Z}(H_{i-1}(BG,\Z),C) \to H^i(BG,C) \to \Hom(H_i(BG,\Z),C) \to 0,
\]

% [原文 p.110]
其中 $H_{i-1}(BG,\Z) = H_{i-1}(G,\Z) = H_{i-1}G$。

此外
\[
  H^1(G,C) \simeq \Hom(H_1G,C)
\]
由此我们推断
\[
  H_1G = G_{(ab)}.
\]
此外
\[
  0 \to \Ext^1_{\Z}(H_1G,C) \to H^2(G,C) \to \Hom(H_2G,C) \to 0.
\]
特别地，若 $G$ 是完美群，即 $G_{(ab)}=0$，则
\[
  H^2(G,C) = \Hom(H_2G,C).
\]
对应于恒等同态 $\id\in\Hom(H_2G,H_1G)$ 的中心扩张如下
\[
  \begin{tikzcd}
    1 \arrow[r] & H_2G \arrow[r, "i"] \arrow[d, "u"'] & \wt G \arrow[r] \arrow[d] & G \arrow[r] \arrow[d, equal] & 1 \\
    1 \arrow[r] & C \arrow[r] & u_*\wt G \arrow[r] & G \arrow[r] & 1
  \end{tikzcd}
\]
其中 $u : H_2G \to C$。

% [Tutor Note: p.110 原文写 $\id\in\Hom(H_2G,H_1G)$（500 dpi 复核无误），但这只有在 $H_1G=H_2G$ 时才有意义；此处应为 $\id\in\Hom(H_2G,H_2G)$——取 $u=\id$ 的推出恰好还原泛中心扩张 $\wt G$。正文照译原文。]
给定任意同态 $u$，有一个推出，它给出 $G$ 被 $C$ 的中心扩张，其中
\[
  u_*\wt G = \frac{C\times\wt G}{\{(-u(x), i(x)) \mid x\in H_2G\}}
\]
且 $i : H_2G \into \wt G$。

因此 $G$ 被 $C$ 的任意中心扩张都由唯一的同态 $u : H_2G \to C$ 诱导。

% [原文 p.111]
\begin{prop}{1}
群 $\wt G$ 是完美的。
\end{prop}

\begin{proof}
有下述交换图
\[
  \begin{tikzcd}
    1 \arrow[r] & H_2G \arrow[r] \arrow[d, "u"'] & \wt G \arrow[r] \arrow[d] & G \arrow[r] \arrow[d, equal] & 1 \\
    1 \arrow[r] & B \arrow[r, "i"] \arrow[d, "i"'] & [\wt G,\wt G] \arrow[r] \arrow[d, "i"'] & G \arrow[r] \arrow[d, equal] & 1 \\
    1 \arrow[r] & H_2G \arrow[r, "i"] & \wt G \arrow[r] & G \arrow[r] & 1
  \end{tikzcd}
\]
其中 $B = H_2G \cap [\wt G,\wt G]$，且 $i$ 表示包含映射。于是由 $u : H_2G \to B$，我们将得到 $i\circ u = 1_{H_2G}$，因为 $i_*\circ u_*(G) = i_*[\wt G,\wt G] = \wt G$。因此 $B = H_2G \Rightarrow [\wt G,\wt G] = \wt G$。
\end{proof}

% [Tutor Note: p.111 原文的 $i_*\circ u_*(G)$ 记号含混（照译），按上下文应读作 $u_*(H_2G)=B$ 且 $i_*$ 把 $B$ 送到 $[\wt G,\wt G]$，从而 $i\circ u=\id_{H_2G}$，于是 $B=H_2G$，即 $H_2G=[\wt G,\wt G]\cap H_2G\subset[\wt G,\wt G]$；又 $[\wt G,\wt G]$ 满射到 $G$ 且核为 $[\wt G,\wt G]\cap H_2G=H_2G=\operatorname{Ker}(\wt G\to G)$，故 $[\wt G,\wt G]=\wt G$。]
\begin{prop}{2}
有
\[
  H_2\wt G = 0,
\]
即 $\wt G$ 没有非平凡的中心扩张，换言之每个中心扩张都分裂。
\end{prop}

\begin{proof}
给定序列
\[
  E \xra{q} \wt G \xra{p} G
\]
其中 $E$ 是 $\wt G$ 的完美中心扩张，且 $p$ 与 $q$ 都是满射，我们断言 $E$ 是 $G$ 的中心扩张。则
\[
  1 \to \operatorname{Ker}q \to \operatorname{Ker}(pq) \to \operatorname{Ker}p \to 1.
\]
$E$ 作用在上述阿贝尔群的正合序列上，它在 $\operatorname{Ker}q$ 与 $\operatorname{Ker}p$ 上的作用平凡，于是我们得到同态 $E \to \Hom(\operatorname{Ker}p,\operatorname{Ker}q)$。因此 $E$ 在 $\operatorname{Ker}(pq)$ 上的作用也平凡。

% [原文 p.112]
于是若 $E = \wt{\wt G}$，则它是 $\wt G$ 的完美中心扩张。有
\[
  \begin{tikzcd}
    1 \arrow[r] & H_2\wt G \arrow[r] & \wt{\wt G} \arrow[r] & \wt G \arrow[r] \arrow[d] & 1 \\
    & & \wt G \arrow[u] & G &
  \end{tikzcd}
\]
而 $\wt G$ 的泛性质蕴涵 $\wt G$ 提升到 $\wt{\wt G}$ 之中。因此
\[
  \wt{\wt G} = \wt G \times H_2\wt G \Rightarrow H_2\wt G = 0.
\]
\end{proof}

\para{例}
\begin{enumerate}[label=\arabic*.]
\item 设 $G = A_5$，即阶为 60 的单非阿贝尔群（它给出二十面体的旋转对称）。$\wt G$ 与 $H_2A_5$ 是什么？

% [Tutor Note: p.112 原文「20-hedron」即 icosahedron（二十面体），注意不是十二面体；$A_5$ 是其旋转对称群，阶为 60。另：$\wt{\wt G}$ 是「$\wt G$ 的泛中心扩张」的记号，不要漏掉外层的波浪号。]
我们首先考虑正合列
\[
  1 \to \{\pm 1\} \to \SU_2 \to \SO_3 \to 1.
\]
回忆拓扑上 $\SU_2 = S^3$，它是 $\C^2$ 的自同构群的子群；而群 $\SO_3$ 是 $\C P^1$ 的自同构群的子群，且拓扑上 $\C P^1 = S^2$。此外 $A_5$ 是 $\SO_3$ 的子群。于是我们有下述图表：
\[
  \begin{tikzcd}
    1 \arrow[r] & \Z_2 \arrow[r] \arrow[d, equal] & S^3 \arrow[r] & \SO_3 \arrow[r] & 1 \arrow[d, equal] \\
    1 \arrow[r] & \Z_2 \arrow[r] & E \arrow[r] & G \arrow[r] & 1
  \end{tikzcd}
\]
我们知道 $E \subset S^3$ 且 $G \subset \SO_3$。我们定义
\[
  M := S^3/E = \SO_3/G
\]
而 $M$ 是一个可定向的 Poincaré 同调 3-球。

% [Tutor Note: p.112 图中两处竖直双线分别连接上下两行的 $\Z_2$ 与最右端的 $1$（不是连接 $\SO_3$ 与 $G$）：核对页图，第二条双线位于两行末端的「$1$」之间。$S^3/E=\SO_3/G$ 成立是因为 $E$ 恰是 $G$ 在二重覆盖 $\SU_2=S^3\to\SO_3$ 下的原像，故 $E/\Z_2=G$，两边商去同一有限群的作用。]
% [原文 p.113]
因为
\[
  \pi_iS^3 = 0
\]
对 $i=0,1,2$，所以 $M$「接近于」成为 $BE$。对 $i=1,2$，我们有
\[
  \pi_i(M) = \pi_iS^3 = 0
\]
由此推得
\[
  H_iM = H_iE.
\]
$M$ 是可定向 3-流形，因此由 Poincaré 对偶我们有 $H_2M = 0$，因为 $H_1M = H_1E = 0$。于是 $H_2E = 0$。最后
\[
  H_1E = H_2E = 0 \Rightarrow E = A_5
\]
以及
\[
  H_1E = H_2E = 0 \Rightarrow H_2A_5 = \Z_2.
\]
\noindent$\square$

% [Tutor Note: p.113 原文作「Poincre duality」（拼写笔误），应为 Poincaré 对偶：$M$ 是可定向闭 3-流形，对偶给出 $H_2M\cong H^1(M;\Z)\cong\Hom(H_1M,\Z)$（$H_0M=\Z$ 自由），故 $H_1M=H_1E=0$ 时 $H_2M=0$。正文照译原文。]
\item 我们将提及 K 理论与 Lie 群同伦群之间的一些类比。我们知道
\[
  S_3 \onto \SO_3 \subset \OO_3
\]
其中 $\pi_0(\OO_3) = \SO_3$ 是 $\OO_3$ 的单位的连通分支，而 $S^3$ 是 $\SO_3$ 的万有覆盖。

% [Tutor Note: p.113 第一个记号原文印作 $S_3$（下标 3），按上下文应为 $S^3$（见同页下一行与 p.112「$S^3$ 是 $\SO_3$ 的泛覆盖」）；$\pi_0(\OO_3)=\SO_3$ 亦有误：应为 $\pi_0(\OO_3)=\Z_2$，而 $\SO_3$ 是 $\OO_3$ 中含单位的连通分支。正文均照译原文。]
如果我们把「连通」一词换成「完美」，我们有下述类比：
\[
  \operatorname{St}A \to EA \subset \GL A
\]
其中 $EA$ 是 $\GL A$ 的最大完美子群，而 $\operatorname{St}A = \wt E A$ 是万有覆盖。

无限正交群 $\OO_\infty = \OO$ 与 K 理论更为接近：
\[
  \operatorname{Spin} \onto \SO \subset \OO
\]
且 $\pi_0(\OO) = K^{-1}O(\mathrm{pt})$ 是包含 $\SO\into\OO$ 的核，而 $\pi_1(\OO) = K^{-2}O(\mathrm{pt})$ 是满射 $\operatorname{Spin}\onto\SO$ 的核。

% [Tutor Note: p.113 原文的 $K^{-1}O(\mathrm{pt})$、$K^{-2}O(\mathrm{pt})$（500 dpi 复核：上标 $-1$、$-2$ 在 $K$ 之后、$O$ 之前）应读作实 K 理论的 $KO^{-1}(\mathrm{pt})$ 与 $KO^{-2}(\mathrm{pt})$：$KO^{-1}(\mathrm{pt})=\pi_0(\OO)=\Z_2$，$KO^{-2}(\mathrm{pt})=\pi_1(\OO)=\Z_2$，与正文所述「核」一致。]
\end{enumerate}
